\subsection{The interpreter lock and the device}
\label{sec:gil}

Python threads overlap only while they are outside the interpreter: a blocking system
call releases the global interpreter lock (GIL) so that another thread can run. In
CPython 3.12.x and 3.13.x, \texttt{fcntl.ioctl} releases it only when it can copy the
argument into a 1{,}024-byte stack buffer. A larger writable buffer is passed to the
kernel in place, and the call runs with the GIL held (\texttt{fcntl\_ioctl\_impl} in
\texttt{Modules/fcntlmodule.c}, unchanged from v3.12.3 to v3.13.9; v3.14.0 releases
the GIL on this path, which we read in the source but did not run). The fork's FIEMAP
request has room for 32 extents, $32 + 32\cdot56 = 1{,}824$ bytes, so every FIEMAP call
holds the lock for as long as the kernel takes, including any device read for an
extent-tree block. The largest request that takes the releasing path holds 17 extents
(984 bytes).

Table~\ref{tab:gil} isolates the effect. Each of 200 files has 16 extents, so FIEMAP
must read one extent-tree leaf per file; before each timed pass we drop the caches and
\texttt{stat} every file, which brings the inodes in through readahead and leaves the
leaf blocks cold. Both request sizes return identical extents for every file, and every
pass reads 1.0 block per file from the device, so one thread takes the device's
latency per file (8.0\,ms under dm-delay, about 6\,ms on NBD; \S\ref{sec:method}).

When the device serves many requests at once (dm-delay, or NBD with 32 server
threads), 16 threads finish the 984-byte requests 14.6$\times$ and 14.9$\times$ faster
than one thread, but the 1{,}824-byte requests, ours or the fork's own
\texttt{FiemapAnalyzer}, gain nothing: 1.00$\times$ on both devices. Sixteen processes
issuing the same 1{,}824-byte requests gain 15.0$\times$ and 14.9$\times$, so the
request size matters only because of the lock. On the device that serves one request at
a time, nothing gains more than 1.09$\times$. The fork's process engine reaches
8.1$\times$ and 7.8$\times$; its 16 workers take about 0.1\,s longer than ours, which we
did not investigate. With no delay, the fork's thread pool is slower than one thread
(0.107\,s against 0.076\,s).

\begin{table}[t]
\caption{Speedup of 16 workers over one for FIEMAP calls that each wait for one device
read, by request size, threads or processes, and device. Median of five passes, 95\%
bootstrap interval.}
\label{tab:gil}
\small
\inputtable{../results/tab_gil.tex}
\end{table}

Two consequences follow. First, the condition the fork needs most, that waits overlap,
fails exactly for the files FragPicker exists to fix: the fragmented ones, whose
extents do not fit in the inode. Second, the user-space delay hid this as well,
because \texttt{time.sleep} releases the GIL and runs before the ioctl. The fork's
process engine would overlap the waits, but \texttt{--parallel} does not use it, and it
would overlap work that the analysis does not need (\S\ref{sec:correct}).
